% Part: normal-modal-logic
% Chapter: completeness
% Section: modalities-ccs

\documentclass[../../../include/open-logic-section]{subfiles}

\begin{document}

\olfileid{nml}{com}{mod}

\olsection{మోడల్ సంయోజకాలు, సంపూర్ణ అవిరుద్ధ సమితులు}

\begin{explain}
ఏదో ఒక నార్మల్ మోడల్ తర్కం~$\Sigma$లోని సంపూర్ణ
$\Sigma$-అవిరుద్ధ సమితులు~$\Delta$యే లోకాలుగా ఉండే
నమూనా~$\mModel{M^\Sigma}$ను నిర్మించినప్పుడు, ఈ ``లోకాల''
మధ్య ప్రాప్యత సంబంధం~$R^\Sigma$ను కూడా నిర్వచించాలి.
ప్రాప్యత సంబంధాన్ని (అలాగే విలువ నిర్ణయం~$V^\Sigma)$ను
$\mSat{M^\Sigma}{!A}[\Delta]$ అయితే, అప్పుడు మాత్రమే
$!A \in \Delta$ అయ్యేలా నిర్వచించాలనుకుంటున్నాం.
దీన్ని ఎలా చేయాలి?

ప్రాప్యత సంబంధాన్ని నిర్వచించిన తరువాత, లోకంలో సత్యం యొక్క
నిర్వచనం ప్రకారం \iftag{prvBox}{$R^\Sigma\Delta\Delta'$
  అయిన ప్రతి $\Delta'$కు $\mSat{M^\Sigma}{!A}[\Delta']$
  అయితే, అప్పుడు మాత్రమే $\mSat{M^\Sigma}{\Box!A}[\Delta]$}{
  $R^\Sigma\Delta\Delta'$ అయ్యే కనీసం ఒక $\Delta'$కు
  $\mSat{M^\Sigma}{!A}[\Delta']$ అయితే, అప్పుడు మాత్రమే
  $\mSat{M^\Sigma}{\Diamond!A}[\Delta]$}.
$\mSat{M^\Sigma}{!A}[\Delta]$ అయితే, అప్పుడు మాత్రమే $!A \in
\Delta$ అని నిరూపించేటప్పుడు, ప్రత్యేకించి మోడల్ సంయోజకంతో
మొదలయ్యే \tetoken{సూత్రాల}{!!{formula}s}కూ ఇది వర్తించాలి;
అంటే, \iftag{prvBox}{$\mSat{M^\Sigma}{\Box!A}[\Delta]$
  అయితే, అప్పుడు మాత్రమే $\Box !A \in \Delta$}{
  $\mSat{M^\Sigma}{\Diamond!A}[\Delta]$ అయితే, అప్పుడు మాత్రమే
  $\Diamond !A \in \Delta$}. ఈ అవసరాన్ని
\iftag{prvBox}{$\Box !A$}{$\Diamond !A$}కు లోకంలో
సత్య నిర్వచనంతో కలిపితే:
\[
\iftag{prvBox}{%
  \Box!A \in \Delta \text{ అయితే, అప్పుడు మాత్రమే } !A \in \Delta' \text{
    అన్నది } R^\Sigma\Delta\Delta' \text{ అయ్యే ప్రతి $\Delta'$కు}}{%
  \Diamond!A \in \Delta \text{ అయితే, అప్పుడు మాత్రమే } !A \in \Delta' \text{ అన్నది
    కనీసం ఒక } \Delta' \text{కు వర్తిస్తుంది; దానికి } R^\Sigma\Delta\Delta'}
\]
\iftag{prvBox}{ఎడమ నుంచి కుడికి దిశను పరిశీలిద్దాం: ఏ $!A$,
  $R^\Sigma\Delta\Delta'$ అయ్యే ఏ $\Delta'$కైనా,
  $\Box !A \in \Delta$ అయితే $!A \in \Delta'$ అని అది చెబుతుంది.
  $\Box!A \in \Delta$ అయ్యే ప్రతి సందర్భంలో $!A \in \Delta'$ అయినప్పుడే
  $R^\Sigma\Delta\Delta'$ అని నిర్ణయిస్తే ఇది నెరవేరుతుంది.
  ఈ ``అయితే, అప్పుడు మాత్రమే'' యొక్క కుడివైపు షరతును
  $\Setabs{!A}{\Box!A \in \Delta} \subseteq \Delta'$
  అని సంక్షిప్తంగా రాయవచ్చు.}{కుడి నుంచి ఎడమకు దిశను
  పరిశీలిద్దాం: ఏ $!A$, $R^\Sigma\Delta\Delta'$ అయ్యే ఏ
  $\Delta'$కైనా, $!A \in \Delta'$ అయితే $\Diamond!A \in
  \Delta$ అని అది చెబుతుంది. $!A \in \Delta'$ అయ్యే ప్రతి సందర్భంలో
  $\Diamond!A \in \Delta$ అయినప్పుడల్లా $R^\Sigma\Delta\Delta'$
  అని నిర్ణయిస్తే ఇది నెరవేరుతుంది. ఈ ``అయితే, అప్పుడు మాత్రమే''
  యొక్క కుడివైపు షరతును $\Setabs{\Diamond!A}{!A \in \Delta'}
  \subseteq \Delta$ అని సంక్షిప్తంగా రాయవచ్చు.}

ప్రశ్న ఏమిటంటే: $R^\Sigma$ యొక్క ఈ నిర్వచనం నిజంగా
\iftag{prvBox}{$\Box !A \in \Delta$ అయితే, అప్పుడు మాత్రమే
  $\mSat{M^\Sigma}{\Box !A}[\Delta]$ అని హామీ ఇస్తుందా?
  \iftag{prvDiamond}{అలాగే }{}}{}\iftag{prvDiamond}{$\Diamond !A
  \in \Delta$ అయితే, అప్పుడు మాత్రమే $\mSat{M^\Sigma}{\Diamond
  !A}[\Delta]$ అని హామీ ఇస్తుందా?}{} రాబోయే కొన్ని ఫలితాలు
దీనిని స్థాపిస్తాయి.
\end{explain}

\begin{defn}
  $\Gamma$ \tetoken{సూత్రాల}{!!{formula}s} సమితి అయితే,
  \begin{align*}
    \Box\Gamma & = \Setabs{\Box !B}{!B \in \Gamma}\\
    \Diamond\Gamma & = \Setabs{\Diamond !B}{!B \in \Gamma}\\
    \intertext{మరియు}
    \Box^{-1}\Gamma & = \Setabs{!B}{\Box !B \in \Gamma}\\
    \Diamond^{-1}\Gamma & = \Setabs{!B}{\Diamond !B \in \Gamma}\\
  \end{align*}
\end{defn}

అంటే, $\Box\Gamma$ అనేది $\Gamma$ నుంచి ఏర్పడి, $\Gamma$లోని
ప్రతి \tetoken{సూత్రం}{!!{formula}}
ముందు $\Box$ చేర్చితే లభిస్తుంది;
$\Gamma$లో మొదట $\Box$ ఉన్న \tetoken{సూత్రాల}{!!{formula}s}
నుంచి ఆ మొదటి $\Box$ను తొలగిస్తే వచ్చే సమితి
$\Box^{-1}\Gamma$. ఈ నిర్వచనం స్వతంత్రంగా అంత
ముఖ్యమైనది కాదు, కానీ సంకేతరచనను బాగా సరళం చేస్తుంది.

$\Box\Box^{-1}\Gamma \subseteq \Gamma$ అని గమనించండి:
\[
\Box\Box^{-1}\Gamma = \Setabs{\Box!B}{\Box!B \in \Gamma}
\]
అంటే, $\Gamma$లో $\Box$తో మొదలయ్యే
\tetoken{సూత్రాల}{!!{formula}s} సమితే ఇది.

\begin{lem}\ollabel{lem:box1}
  $\Gamma \Proves[\Sigma] !A$ అయితే $\Box\Gamma
  \Proves[\Sigma] \Box !A$.
\end{lem}

\begin{proof}
  $\Gamma \Proves[\Sigma] !A$ అయితే, $!B_1$, \dots, $!B_k
  \in \Gamma$ ఉండి, $\Sigma \Proves !B_1 \lif (!B_2 \lif
  \cdots (!B_k \lif !A)\cdots)$. $\Sigma$ నార్మల్ కనుక
  \RK{} నియమం ప్రకారం $\Sigma \Proves \Box!B_1 \lif
  (\Box!B_2 \lif \cdots (\Box!B_k \lif \Box!A)\cdots)$.
  స్పష్టంగా $\Box!B_1$, \dots, $\Box!B_k \in \Box\Gamma$;
  కనుక నిర్వచనం ప్రకారం $\Box\Gamma \Proves[\Sigma] \Box !A$.
  \sourcecorrection{OLTENMLCOMMOD-001}{స్థిర మూలంలో $!B_1$ నుంచి
    $!B_k$ వరకు సాక్షులను పేర్కొని, రెండు అంతర్నిహితార్థ
    శ్రేణుల్లో చివరి సూచికను $n$గా రాసింది. అదే సాక్ష్యాల
    శ్రేణి కావడానికి ఆ రెండు చోట్ల $k$ అని సరిచేశాం.}
\end{proof}

\begin{lem}\ollabel{lem:box2}
  $\Box^{-1} \Gamma \Proves[\Sigma] !A$ అయితే $\Gamma
  \Proves[\Sigma] \Box!A$.
\end{lem}

\begin{proof}
  $\Box^{-1}\Gamma \Proves[\Sigma] !A$ అనుకుందాం; అప్పుడు
  \olref{lem:box1} ప్రకారం $\Box\Box^{-1}\Gamma
  \Proves[\Sigma] \Box!A$. అయితే $\Box\Box^{-1}\Gamma
  \subseteq \Gamma$ కనుక ఏకదిశత్వం వల్ల $\Gamma
  \Proves[\Sigma] \Box!A$ కూడా.
  \sourcecorrection{OLTENMLCOMMOD-002}{స్థిర మూలంలోని మధ్యంతర
    వ్యుత్పాద్యతకు $\Sigma$ సూచిక లేదు. సూచించిన మొదటి ఉపసిద్ధాంతం
    $\Sigma$-సాపేక్ష వ్యుత్పాద్యతనే ఇస్తుంది; అదే సూచికను
    ఇక్కడ స్పష్టంగా చేర్చాం.}
\end{proof}

\iftag{prvBox}{%
% Box is primitive

\begin{prop}\ollabel{prop:box}
  $\Gamma$ సంపూర్ణ $\Sigma$-అవిరుద్ధమైతే, ప్రతి
  సంపూర్ణ $\Sigma$-అవిరుద్ధ $\Delta$కు $\Box^{-1} \Gamma
  \subseteq \Delta$ అయినప్పుడు $!A \in \Delta$ అవుతుంది
  అయితే, అప్పుడు మాత్రమే $\Box !A \in \Gamma$.
\end{prop}

\begin{proof}
  $\Gamma$ సంపూర్ణ $\Sigma$-అవిరుద్ధం అనుకుందాం. ``అయితేనే''
  దిశ సులభం: $\Box !A \in \Gamma$ మరియు $\Box^{-1}\Gamma
  \subseteq \Delta$ అనుకుందాం. $\Box!A \in \Gamma$ కనుక
  $!A \in \Box^{-1}\Gamma \subseteq \Delta$; అందువల్ల
  $!A \in \Delta$.

  ``అయితే'' దిశకు విపర్యయాన్ని నిరూపిద్దాం: $\Box!A \notin
  \Gamma$ అనుకుందాం. $\Gamma$ సంపూర్ణ $\Sigma$-అవిరుద్ధం
  కనుక నిగమన పరంగా సంవృతం; అందువల్ల $\Gamma
  \Proves/[\Sigma] \Box !A$. \olref{lem:box2} ప్రకారం
  $\Box^{-1}\Gamma \Proves/[\Sigma] !A$.
  \olref[prf][con]{prop:consistencyfacts}\olref[prf][con]{prop:consistencyfacts-b}
  ప్రకారం $\Box^{-1}\Gamma \cup \{ \lnot!A \}$
  $\Sigma$-అవిరుద్ధం. లిండెన్‌బామ్ ఉపసిద్ధాంతం ప్రకారం
  $\Box^{-1}\Gamma \cup \{ \lnot!A \} \subseteq \Delta$
  అయ్యే సంపూర్ణ $\Sigma$-అవిరుద్ధ సమితి~$\Delta$ ఉంది.
  అవిరుధ్యం వల్ల $!A \notin \Delta$.
\end{proof}
}{%
% Box is not primitive, only Diamond is

\begin{prop}\ollabel{prop:diamond}
  $\Gamma$ సంపూర్ణ $\Sigma$-అవిరుద్ధమైతే,
  $\Diamond !A \in \Gamma$ అయితే, అప్పుడు మాత్రమే
  $\Diamond\Delta \subseteq \Gamma$ మరియు $!A \in \Delta$
  అయ్యే సంపూర్ణ $\Sigma$-అవిరుద్ధ సమితి $\Delta$ ఒకటి ఉంటుంది.
\end{prop}

\begin{proof}
  $\Gamma$ సంపూర్ణ $\Sigma$-అవిరుద్ధం అనుకుందాం. కుడి
  నుంచి ఎడమకు దిశ సులభం: $\Diamond\Delta \subseteq \Gamma$
  మరియు $!A \in \Delta$ అయ్యే సంపూర్ణ $\Sigma$-అవిరుద్ధ
  $\Delta$ను తీసుకుందాం. అప్పుడు $\Diamond!A \in
  \Diamond\Delta \subseteq \Gamma$; అంటే $\Diamond!A \in
  \Gamma$.

  ఎడమ నుంచి కుడికి దిశకు $\Diamond !A \in \Gamma$
  అనుకుందాం. $!A \in \Delta$, $\Diamond\Delta \subseteq
  \Gamma$ అయ్యే సంపూర్ణ $\Sigma$-అవిరుద్ధ $\Delta$ ఉందని
  చూపాలి. $\Diamond !A \in \Gamma$, $\Gamma$
  $\Sigma$-అవిరుద్ధం కనుక $\Box\lnot !A \notin \Gamma$.
  $\Gamma$ సంపూర్ణ $\Sigma$-అవిరుద్ధం కనుక నిగమన పరంగా
  సంవృతం; అందుచేత $\Gamma \Proves/[\Sigma] \Box\lnot !A$.
  \olref{lem:box2} ప్రకారం $\Box^{-1}\Gamma
  \Proves/[\Sigma] \lnot !A$.
  \olref[prf][con]{prop:consistencyfacts}\olref[prf][con]{prop:consistencyfacts-b}
  ప్రకారం $\Box^{-1}\Gamma \cup \{ !A \}$
  $\Sigma$-అవిరుద్ధం. లిండెన్‌బామ్ ఉపసిద్ధాంతం ప్రకారం
  $\Box^{-1}\Gamma \cup \{ !A \} \subseteq \Delta$
  అయ్యే సంపూర్ణ $\Sigma$-అవిరుద్ధ సమితి~$\Delta$ ఒకటి
  ఉంటుంది. స్పష్టంగా $!A \in \Delta$. ఇప్పుడు
  $\Diamond\Delta \subseteq \Gamma$ అని చూడాలి. అలా
  కాకపోతే, ఏదో $!B \in \Delta$కు $\Diamond !B \notin
  \Gamma$. అప్పుడు $\Gamma$ సంపూర్ణ $\Sigma$-అవిరుద్ధం
  కనుక $\Box\lnot !B \in \Gamma$. కానీ $\lnot !B \in
  \Box^{-1}\Gamma \subseteq \Delta$ కూడా; ఇది $\Delta$ను
  విరుద్ధంగా చేస్తుంది.
\end{proof}
}

\begin{lem}\ollabel{lem:box-iff-diamond}
  $\Gamma$, $\Delta$ సంపూర్ణ $\Sigma$-అవిరుద్ధ సమితులు
  అనుకుందాం. అప్పుడు $\Box^{-1}\Gamma \subseteq \Delta$
  అయితే, అప్పుడు మాత్రమే $\Diamond\Delta \subseteq \Gamma$.
\end{lem}

\begin{proof}
  ``అయితేనే'' దిశ: $\Box^{-1}\Gamma \subseteq \Delta$
  అనుకుని, $\Diamond!A \in \Diamond\Delta$ (అంటే $!A \in
  \Delta$) అనుకుందాం. $\Diamond!A \in \Gamma$ అని చూపడానికి
  $\Box\lnot!A \notin \Gamma$ అని చూపడం చాలు; ఎందుకంటే అప్పుడు
  గరిష్ఠత వల్ల $\lnot\Box\lnot !A \in \Gamma$. ఇప్పుడు
  $\Box\lnot!A \in \Gamma$ అయితే పరికల్పన వల్ల $\lnot!A
  \in \Delta$; ఇది $!A \in \Delta$ కూడా ఉన్న $\Delta$
  అవిరుధ్యానికి విరుద్ధం. కనుక కావలసినట్లే $\Box\lnot!A
  \notin \Gamma$.

  ``అయితే'' దిశ: $\Diamond\Delta \subseteq \Gamma$
  అనుకుందాం. విపర్యయంగా వాదిద్దాం: $\Box!A \notin \Gamma$
  అని చూపడానికి $!A \notin \Delta$ అనుకుందాం.
  $!A \notin \Delta$ అయితే గరిష్ఠత వల్ల $\lnot!A \in
  \Delta$; కాబట్టి పరికల్పన వల్ల $\Diamond\lnot!A \in \Gamma$.
  అయితే నార్మల్ మోడల్ తర్కంలో $\Diamond\lnot!A$,
  $\lnot\Box !A$కు తుల్యం. తరువాతిది $\Gamma$లో ఉంటే,
  అవిరుధ్యం వల్ల కావలసినట్లే $\Box!A \notin\Gamma$.
\end{proof}

\iftag{prvBox}{%
  \iftag{prvDiamond}{%
% Box and Diamond are both primitive; prove
% prop:diamond using lem:box-iff-diamond

\begin{prop}\ollabel{prop:diamond}
  $\Gamma$ సంపూర్ణ $\Sigma$-అవిరుద్ధమైతే,
  $\Diamond !A \in \Gamma$ అయితే, అప్పుడు మాత్రమే
  $\Diamond\Delta \subseteq \Gamma$, $!A \in \Delta$
  అయ్యే సంపూర్ణ $\Sigma$-అవిరుద్ధ సమితి $\Delta$ ఒకటి ఉంటుంది.
\end{prop}

\begin{proof}
$\Gamma$ సంపూర్ణ $\Sigma$-అవిరుద్ధం అనుకుందాం. \Dual{}
మరియు సంవృతత వల్ల $\Diamond!A \in \Gamma$ అయితే,
అప్పుడు మాత్రమే $\lnot\Box\lnot !A \in \Gamma$.
$\Gamma$ సంపూర్ణ $\Sigma$-అవిరుద్ధం కనుక
\olref[ccs]{prop:ccs-properties}\olref[ccs]{prop:ccs-lnot}
ప్రకారం $\lnot\Box\lnot !A \in \Gamma$ అయితే, అప్పుడు
మాత్రమే $\Box\lnot !A \notin \Gamma$. \olref{prop:box}
ప్రకారం $\Box\lnot !A \notin \Gamma$ అయితే, అప్పుడు
మాత్రమే $\Box^{-1}\Gamma \subseteq \Delta$ అయ్యే ఏదో
సంపూర్ణ $\Sigma$-అవిరుద్ధ $\Delta$కు $\lnot!A \notin
\Delta$. అటువంటి~$\Delta$ను తీసుకుందాం.
\olref{lem:box-iff-diamond} ప్రకారం $\Box^{-1}\Gamma
\subseteq \Delta$ అయితే, అప్పుడు మాత్రమే $\Diamond\Delta
\subseteq \Gamma$. అలాగే
\olref[ccs]{prop:ccs-properties}\olref[ccs]{prop:ccs-lnot}
ప్రకారం $\lnot !A \notin \Delta$ అయితే, అప్పుడు మాత్రమే
$!A \in \Delta$. కనుక $\Diamond!A \in \Gamma$ అయితే,
అప్పుడు మాత్రమే $\Diamond\Delta \subseteq \Gamma$,
$!A \in \Delta$ అయ్యే ఏదో సంపూర్ణ $\Sigma$-అవిరుద్ధ
$\Delta$ ఉంటుంది.
\end{proof}

}{}}{}

\begin{prob}
  $\Gamma$ సంపూర్ణ $\Sigma$-అవిరుద్ధమైతే,
  \iftag{prvBox}{$\Diamond !A \in \Gamma$ అయితే, అప్పుడు
    మాత్రమే $\Box^{-1}\Gamma \subseteq \Delta$, $!A \in
    \Delta$ అయ్యే సంపూర్ణ $\Sigma$-అవిరుద్ధ సమితి $\Delta$
    ఒకటి ఉంటుందని.}{$\Box !A \in \Gamma$ అయితే, అప్పుడు
    మాత్రమే $\Box^{-1} \Gamma \subseteq \Delta$ అయ్యే
    ప్రతి సంపూర్ణ $\Sigma$-అవిరుద్ధ $\Delta$కు $!A \in
    \Delta$ అవుతుందని.} నిరూపించండి. \emph{ఇందుకోసం
    \olref[nml][com][mod]{lem:box-iff-diamond} ఉపయోగించవద్దు}.
\end{prob}

\end{document}
